\documentclass[10pt,a4paper]{amsart}
\usepackage[margin=19mm]{geometry}
\usepackage{amsmath,amssymb,mathtools}
\usepackage[scr=rsfs,cal=euler]{mathalfa}
\usepackage{xfrac}
\usepackage[unicode,hidelinks,bookmarks=false]{hyperref}
\newcommand{\ie}{i.e.\ }
\newcommand{\cf}{cf.\ }
\newcommand{\resp}{respectively\ }
\newcommand{\Subsection}{\subsection}
\providecommand{\nobreakdash}{\nobreak\hbox{-}}
\setlength{\emergencystretch}{2em}
\multlinegap0pt
\usepackage{bm}
\usepackage{bbm}
\usepackage{dsfont}
\usepackage{xspace}
\usepackage{cancel}
\usepackage{mathtools}
\usepackage{amsthm}
\makeatletter
\@ifundefined{lemma}{\newtheorem{lemma}{Lemma}}{}
\makeatother
\title[On solvmanifolds with complex commutator and constant holomo]{On solvmanifolds with complex commutator and constant holomorphic sectional curvature}
\author{Xin Huang, Fangyang Zheng}
\date{Source: arXiv:2501.00810v1 (CC BY 4.0)}
\begin{document}
\maketitle

\noindent\emph{Source and licence.} On solvmanifolds with complex commutator and constant holomorphic sectional curvature. Version arXiv:2501.00810v1 (CC BY 4.0), distributed under the Creative Commons Attribution 4.0 International licence, as recorded in the arXiv licence metadata for this exact version and shown on the abstract page https://arxiv.org/abs/2501.00810v1. Full rights notice: licenses/arxiv-2501.00810-CC-BY-4.0.txt.\par\smallskip

\noindent\textit{Editorial scope.} Counted unit: the Lemma whose source label is \texttt{lemmagoal} in the article (source0.tex lines 263--268, source label \texttt{lemmagoal}), delivered together with the article's own complete original proof of it (source0.tex lines 270--318, one \texttt{proof} environment of 49 lines). No mathematics is written, repaired or re-derived: statement and proof are the article's own text line for line. Required context retained uncounted: eq:res1 (lines 234--236); eq:res2 (lines 238--240); eq:alpha (lines 249--251); eq:i (lines 253--255); eq:r (lines 257--259). Every definition, display and label of those lines is retained unchanged. Omitted from this package and not redistributed: 1--233, 237--237, 241--248, 252--252, 256--256, 260--262, 269--269, 319--777. Nothing inside a retained excerpt is omitted or altered: every retained body is delivered exactly as the article's lines are written, so no omission receipt of any kind accompanies this package. Citations: the retained excerpts carry no citation command at all, so no bibliography entry of the article is transcribed and no reference is redistributed. Reference adaptation: none was needed and none was made; every reference inside the delivered unit resolves inside it, so no unresolved reference or citation is produced. Printed slips kept literal: no printed word, symbol, hypothesis or number of the counted statement or of its proof was altered. Layout and class: the article's own class is \texttt{amsart} with packages including amsmath, amscd, amssymb, amsthm, bbm, latexsym, amsfonts, graphicx, xy, hyperref, geometry; the wrapper is the lane's pinned \texttt{amsart} editorial wrapper at 10pt on A4, which restates the theorem-like environments the retained text needs and adds the article's own macro definitions that the retained text needs (none), with extra wrapper packages (bm, bbm, dsfont, xspace, cancel, mathtools). The delivered numbering is the unit's own. Assets: no retained span contains a figure environment, a graphics-inclusion command, a PDF-page inclusion, a table or a TikZ picture; no image file is copied into this package, the package declares no asset dependency and no image file is redistributed with it. Licence: version arXiv:2501.00810v1 of this article is distributed under the Creative Commons Attribution 4.0 International licence, as recorded in the arXiv licence metadata for this exact version and shown on the abstract page https://arxiv.org/abs/2501.00810v1. A full rights notice is delivered at licenses/arxiv-2501.00810-CC-BY-4.0.txt. Characters: every retained excerpt is pure ASCII, so the delivered engine, XeLaTeX with its default Latin Modern text font, reports no missing character.
\begin{equation}  \label{eq:res1}
C^{\alpha}_{ab}=D^{a}_{\alpha b}=0, \ \ \ \ \forall \ r\!+\!1\leq \alpha \leq n, \ \forall \ 1\leq a,b\leq n. 
\end{equation}

\begin{equation}  \label{eq:res2}
C^j_{ik}=0  \ \  \mbox{unless} \ j>i \ \mbox{or} \ j>k; \ \ \ \ D^j_{ik}=0 \  \ \mbox{unless} \ i>j, \ \ \ \forall \ 1\leq i,j,k\leq r.
\end{equation}

\begin{equation} \label{eq:alpha}
 c=R_{\alpha \bar{\alpha} \alpha \bar{\alpha} } = - \sum_{s=1}^n |D^{\alpha }_{s\alpha}|^2 \ \leq \,0. 
 \end{equation}

\begin{equation}  \label{eq:i}
 c=R_{i \bar{i} i \bar{i} } =  \sum_{s=1}^n |D^s_{ii}|^2 - \sum_{s=1}^r |D^{i}_{si}|^2 ,
 \end{equation}

\begin{equation}  \label{eq:r}
 c=R_{r \bar{r} r \bar{r} } =  \sum_{s=1}^n |D^s_{rr}|^2 \ \geq \, 0. 
 \end{equation}

\begin{lemma} \label{lemmagoal}
Let $({\mathfrak g},J,g)$ be a Lie algebra with Hermitian structure such that ${\mathfrak g}$ is solvable and  $J{\mathfrak g}'={\mathfrak g}'$. Assume that the holomorphic sectional curvature $H$ of the Chern connection of $g$ is zero. Then under any admissible frame $e$ of ${\mathfrak g}$ it holds
\begin{equation}  \label{eq:goal}
D^{\alpha}_{j\beta}=D^j_{ik}=D^{\alpha}_{ij}=0,  \ \ \ \ \ \ \ \ \forall\ 1\leq i,j,k\leq r, \ \ \ \forall \ r\!+\!1\leq \alpha , \beta \leq n. 
 \end{equation}
\end{lemma}

\begin{proof}
By formula (\ref{eq:alpha}) and the vanishing of the holomorphic sectional curvature, we know that $D^{\alpha}_{j\alpha}=0$ for any $j$ and any $\alpha$. If we make any unitary change on $\{ e_{r\!+\!1},\ldots , e_n\}$, formula (\ref{eq:res1}) and (\ref{eq:res2}) would not be affected and $e$ would still be an admissible frame. So formula (\ref{eq:alpha}) actually gives us $D^{X}_{jX}=0$ for any $j$ and any $X=\sum_{\alpha =r\!+\!1}^n X_{\alpha}e_{\alpha}$. That is, $\sum_{\alpha ,\beta=r\!+\!1}^n \overline{X}_{\alpha} X_{\beta} D^{\alpha}_{j\beta }=0$. This leads to $D^{\alpha}_{j\beta}=0$ for any $j$ and any $\alpha$, $\beta$. 

Analogously, the vanishing of $H$ and (\ref{eq:r}) give us $D^{\ast}_{rr}=0$. Our next goal is to prove the remaining part of (\ref{eq:goal}), namely, 
$D^j_{ik}=D^{\alpha}_{ij}=0$ for any $ i$,\,$j$,\,$k$ and any $ \alpha$. To prove this, let us fix any $1\leq i<k\leq r$ and consider
\begin{eqnarray}
0 & = & 4\widehat{R}_{i\bar{i}k\bar{k}} \ = \ R_{i\bar{i}k\bar{k}} + R_{k\bar{k}i\bar{i}}  + R_{i\bar{k}k\bar{i}} +R_{k\bar{i}i\bar{k}} \nonumber \\
& = & \sum_{s=1}^n \{ \big( |D^s_{ki}|^2 -  |D^k_{si}|^2\big) + \big(|D^s_{ik}|^2 -  |D^i_{sk}|^2\big)  +2\mbox{Re}  \big( D^s_{ki}\overline{D^s_{ik}} - D^i_{si}\overline{D^k_{sk}} - \overline{D^i_{sk}} D^i_{ks} \big) \} \nonumber \\
& = & \sum_{s=1}^n |D^s_{ki}+D^s_{ik}|^2 -  \sum_{s=1}^r \big( |D^k_{si}|^2 + |D^i_{sk}|^2 + 2\mbox{Re} \{ D^i_{si}\overline{D^k_{sk}} + \overline{D^i_{sk}} D^i_{ks} \} \big). \label{eq:iikk}
\end{eqnarray}
In the last equality above, we used the fact that $D^{\ast}_{\alpha \ast}=0$ for any $\alpha >r$ by (\ref{eq:res1}). Let $k=r$ in (\ref{eq:iikk}), and by (\ref{eq:res2}) we get
\begin{equation}  \label{eq:iirr}
 \sum_{s=1}^n |D^s_{ri}+D^s_{ir}|^2 \, = \, \sum_{s=1}^r \big(  |D^i_{sr}|^2 + 2\mbox{Re} \{  \overline{D^i_{sr}} D^i_{rs} \} \big) , \ \ \ \ \ \ \forall \ 1\leq i<r.
 \end{equation}
We claim that the above will lead to the following:
\begin{equation}  \label{eq:claimr}
 D^s_{ri}+D^s_{ir}=D^i_{jr}=D^i_{rj}=0, \ \ \ \ \ \ \forall \ 1\leq i,j\leq r, \ \ \forall \ 1\leq s\leq n.
 \end{equation}
To prove the claim, let us perform induction on $i$. It holds for $i=r$ by (\ref{eq:res2}) and the fact $D^{\ast}_{rr}=0$. Let $i=r\!-\!1$ in (\ref{eq:iirr}). The right-hand side of (\ref{eq:iirr}) vanishes as $s$ has to be $r$ by (\ref{eq:res2}) while $D^{\ast}_{rr}=0$. Therefore we get $D^s_{r,r\!-\!1}+D^s_{r\!-\!1,r} =0$ for any $s$, hence (\ref{eq:claimr}) holds for $i=r\!-\!1$. Next let us take $i=r\!-\!2$ in (\ref{eq:iirr}). On the right hand side, $s$ has to be $r\!-\!1$. But by (\ref{eq:claimr}) for $i=r\!-\!1$ we have $D^{r\!-\!2}_{r\!-\!1,r} + D^{r\!-\!2}_{r,r\!-\!1}=0$. Therefore the right-hand side of (\ref{eq:iirr}) becomes $-| D^{r\!-\!2}_{r\!-\!1,r} |^2$, so we get (\ref{eq:claimr}) for the $i=r\!-\!2$ case. Repeating this process, the Claim (\ref{eq:claimr}) is proved. In particular, we have shown that
\begin{equation*}  \label{eq:claimDr}
 D^i_{jr}=D^i_{rj}=D^r_{ij} =0, \ \ \ \ \ \ \forall \ 1\leq i,j\leq r.
 \end{equation*}
If we look at the $i=r\!-\!1$ case of (\ref{eq:i}), then the second term on the right-hand side vanishes, so we have
\begin{equation}  \label{eq:Dr-1r-1}
 D^s_{r\!-\!1,r\!-\!1}=0, \ \ \ \ \ \ \forall \ 1\leq s\leq n.
 \end{equation}
Now let us take $k=r\!-\!1$ in (\ref{eq:iikk}), and get
\begin{equation}  \label{eq:iir-1r-1}
 \sum_{s=1}^n |D^s_{r\!-\!1,i}+D^s_{i,r\!-\!1}|^2 \, = \, \sum_{s=1}^r \big(  |D^i_{s,r\!-\!1}|^2 + 2\mbox{Re} \{  \overline{D^i_{s,r\!-\!1}} D^i_{r\!-\!1,s} \} \big) , \ \ \ \ \ \ \forall \ 1\leq i<r\!-\!1.
 \end{equation}
It leads to the following by arguing inductively in $i$ in decreasing order as we have seen before:
\begin{equation}  \label{eq:claimr-1}
 D^s_{r\!-\!1,i}+D^s_{i,r\!-\!1}=D^i_{j,r\!-\!1}=D^i_{r\!-\!1,j}=0, \ \ \ \ \ \ \forall \ 1\leq i,j\leq r, \ \ \forall \ 1\leq s\leq n.
 \end{equation}
In particular, we have shown that
$$ D^i_{j,r\!-\!1}=D^i_{r\!-\!1,j}=D^{r\!-\!1}_{ij} =0, \ \ \ \ \ \ \forall \ 1\leq i,j\leq r. $$
In other words, for $D^j_{ik}$ with $1\leq i,j,k\leq r$, it vanishes if any of the index is $r$ or $r\!-\!1$. Repeating this process, we conclude that 
$D^j_{ik}=0$ for all  $1\leq i,j,k\leq r$. This completes the proof of the second equality of (\ref{eq:goal}). To prove the third equality of (\ref{eq:goal}), note that the second term on the right-hand side of (\ref{eq:i}) now vanishes, so it gives us $D^{\alpha}_{ii}=0$. If we perform a unitary change on $\{ e_1, \ldots , e_r\}$, we still have $D^j_{ik}=0$ for all $1\leq i,j,k\leq r$. So what (\ref{eq:i}) gives us is actually $D^{\alpha}_{vv}=0$ for any $v=v_1e_1+\cdots + v_re_r$. In particular we have 
\begin{equation}  \label{eq:Dalphaij}
 D^{\alpha}_{ij}+ D^{\alpha}_{ji} =0 , \ \ \ \ \ \ \ \ \forall \ 1\leq i,j\leq r, \ \ \forall \ r\!+\!1\leq \alpha \leq n.
 \end{equation}
For any fixed $1\leq i\leq r$ and $r\!+\!1\leq \alpha \leq n$, by (\ref{eq:res2}) and the first equality of (\ref{eq:goal}) we have
\begin{eqnarray*}
0 & = & 4\widehat{R}_{i\bar{i}\alpha \bar{\alpha}} \ = \ R_{i\bar{i}\alpha \bar{\alpha}} + R_{\alpha \bar{\alpha}i\bar{i}}  + R_{i\bar{\alpha}\alpha \bar{i}} +R_{\alpha \bar{i}i\bar{\alpha}} \nonumber \\
& = & \sum_{s=1}^r  \{  |D^s_{i\alpha}|^2 -  |D^i_{s\alpha}|^2- |D^{\alpha}_{si}|^2 -  2\mbox{Re}  \big( D^{\alpha}_{si}\overline{D^{\alpha}_{is}}  \big) \} \nonumber \\
& = & \sum_{s=1}^r  \{  |D^s_{i\alpha}|^2 -  |D^i_{s\alpha}|^2 + |D^{\alpha}_{si}|^2  \}.
\end{eqnarray*}
Here in the last equality we used the (\ref{eq:Dalphaij}) for the last term. If we sum $i$ from $1$ to $r$, then the first two terms on the right-hand side of the last line above will cancel each other, and it gives us $D^{\alpha}_{si}=0$ for any $1\leq s,i\leq r$. This completes the proof of the third equality of (\ref{eq:goal}) hence the lemma. 
\end{proof}

\end{document}
